\documentclass[manuscript, screen]{jair}

\setcopyright{none}
\acmDOI{}
\JAIRAE{}
\JAIRTrack{}
\acmVolume{}
\acmArticle{}
\acmMonth{11}
\acmYear{2026}

\RequirePackage[
  datamodel=acmdatamodel,
  style=acmauthoryear,
  backend=biber,
  giveninits=true,
  uniquename=init
  ]{biblatex}
\renewcommand*{\bibopenbracket}{(}
\renewcommand*{\bibclosebracket}{)}
\addbibresource{references.bib}

\usepackage{tabularx}

\newcommand{\lean}[1]{\texttt{\detokenize{#1}}}
\newcommand{\tnI}{\mathit{tn}_I}
\newcommand{\sI}{s_I}
\newcommand{\tn}{\mathit{tn}}
\newcommand{\pre}{\mathit{pre}}
\newcommand{\eff}{\mathit{eff}}
\newcommand{\FM}{\textsc{fond}\textsubscript{FM}}
\newcommand{\MP}{\textsc{fond}\textsubscript{MP}}
\newcommand{\CBa}{\textcite{chen:bercher:icaps-21}}
\newcommand{\CBb}{\textcite{chen:bercher:icaps-22}}
\newcommand{\DB}{\textcite{dekker:behnke:icaps-25}}
\newcolumntype{L}{>{\raggedright\arraybackslash}X}

\newtheorem{deviation}{Deviation}

\begin{document}

\title[FOND HTN Planning in Lean]{Where the Lean Formalisation of FOND HTN Planning Differs
  from the Literature}

\author{Gregor Behnke}
\email{g.behnke@uva.nl}
\affiliation{%
  \institution{University of Amsterdam}
  \city{Amsterdam}
  \country{The Netherlands}}

\renewcommand{\shortauthors}{Behnke}

\begin{abstract}
  We formalise fully observable nondeterministic (FOND) HTN planning in Lean~4 with Mathlib
  \parencite{moura:ullrich:cade-21,mathlib:cpp-20}, covering the fixed-method solution criteria
  of \CBa{} (weak and strong, linearisation- and outcome-dependent), the solvable belief states
  of \DB{}, and the method-based policies of \CBb{} (weak, strong cyclic and strong). This report
  lists every place where the formal definitions deviate from the papers and why, and documents
  the formally proved relationships between the notions together with the deviations their
  proofs required. All theorems depend only on the axioms \lean{propext},
  \lean{Classical.choice} and \lean{Quot.sound}.
\end{abstract}

\maketitle

\section{Overview}

The development lives in the Lean library \lean{FondHtn} and builds on the deterministic HTN
library \lean{HtnPlanning}, from which it takes task networks, methods, isomorphism of task
networks, decomposition, progression and the structural problem classes. Its modules are:

\begin{center}
\begin{tabularx}{\linewidth}{@{}lL@{}}
  \toprule
  module & content \\
  \midrule
  \lean{FondHtn.Basic} & nondeterministic actions, domains, problems, problem classes \\
  \lean{FondHtn.FixedMethods} & \CBa{}: LD and OD solutions; \DB{}: solvable belief states \\
  \lean{FondHtn.FixedMethodsLemmas} & strong OD solutions $\Leftrightarrow$ solvable belief states \\
  \lean{FondHtn.MethodPolicies} & \CBb{}: method-based policies and their solution criteria \\
  \bottomrule
\end{tabularx}
\end{center}

\noindent
We refer to the three papers as [CB21], [CB22] and [DB25]. Section~\ref{sec:model} lists
deviations in the common model, Section~\ref{sec:fm} those of the fixed-method criteria,
Section~\ref{sec:mp} those of the method-based criteria, and Section~\ref{sec:results} the
proved results.

\section{The Common Model}\label{sec:model}

\begin{deviation}[Negative preconditions]\label{dev:neg}
  An action is a triple $(\Pi^+,\Pi^-,E)$ of a positive precondition, a negative precondition and
  a set $E$ of effects $(\mathit{add},\mathit{del})$ (\lean{FOND.Action}), as in [DB25]. [CB21]
  and [CB22] have positive preconditions only; their actions are the special case $\Pi^-=\emptyset$.
  \emph{Reason:} the hardness constructions of [DB25] and of the planned work use negative
  preconditions; the generalisation costs nothing.
\end{deviation}

\begin{deviation}[Well-formed actions]\label{dev:wf}
  The papers implicitly assume that every action has finitely many and at least one effect
  ($1 \le |\eff(a)| < \infty$). This is the separate predicate \lean{FOND.Action.WF} and part of
  \lean{FOND.Problem.WF}; it is not built into the type. \emph{Consequence:} for an action without
  effects, $\gamma(a,s)=\emptyset$. Belief-state solvability and strong OD solutions then hold
  vacuously after executing $a$, whereas in the execution structure of [CB22] the node becomes
  terminal (and hence fails unless it is a goal node). Under \lean{WF} the two readings agree.
\end{deviation}

\begin{deviation}[Types instead of finite sets]
  Facts, primitive and compound task names and task identifiers are types \lean{F}, \lean{A},
  \lean{C}, \lean{I}; states are arbitrary subsets of \lean{F}. Finiteness is to be imposed by
  finite descriptions, as in \lean{HtnPlanning}. The action mapping $\delta$ is total on
  \lean{A}, which plays the role of $N_P$.
\end{deviation}

\begin{deviation}[Fresh identifiers]\label{dev:fresh}
  Decomposition (\lean{HTN.DecompStep}) inserts a well-formed copy of the method's network with
  identifiers disjoint from the decomposed network. This exists only if enough identifiers are
  available, so the notions involving decomposition are intended for an infinite type \lean{I};
  theorems that need it assume \lean{[Infinite I]}.
\end{deviation}

\begin{deviation}[Problem classes]\label{dev:classes}
  The classes primitive, totally ordered, regular, acyclic and tail-recursive
  ([CB21, Defs.~2.14, 3.2--3.6], [CB22, Defs.~8--9]) depend only on $\tnI$ and the methods. They
  are inherited from \lean{HtnPlanning} by applying its predicates to the \emph{skeleton}
  (\lean{FOND.Problem.skeleton}), the deterministic problem with the same $\tnI$ and methods and
  trivial actions. In \lean{HtnPlanning} the conditions on methods (total order, regularity, the
  stratification conditions) are imposed only on methods \emph{reachable} from $\tnI$, whereas the
  papers impose them on all methods. Unreachable methods can never be applied, so this does not
  affect any solution notion; it only makes the classes slightly larger.
\end{deviation}

\section{Fixed Methods}\label{sec:fm}

Achievable primitive task networks [CB21, Def.~2.5] and linearisation-dependent (LD) solutions
[CB21, Def.~2.7] are formalised literally (\lean{AchievablePrimitive},
\lean{IsWeakLDSolution}, \lean{IsStrongLDSolution}). The sets $S_i$ of Def.~2.7 are threaded
through the recursion \lean{StrongExec}; weak executability guesses the outcomes
(\lean{WeakExec}).

\begin{deviation}[Outcome-dependent solutions]\label{dev:od}
  Taken literally, [CB21, Defs.~2.9--2.11] do not express the intended notion:
  \begin{itemize}
    \item The execution structure $K=\langle Q,R\rangle$ of Def.~2.9 adds a state $s$ to $Q$
      whenever $s$ lies in \emph{some} set executable by \emph{some} task set $T'$ with
      $\pi(T',s)$ defined. The edges thus forget the history $T'$ along which $s$ was reached,
      so an execution sequence may follow edges created for a different history.
    \item By Def.~2.10 every prefix of an execution sequence is again an execution sequence. Hence
      ``every execution sequence of $K$ is complete'' (strong OD, Def.~2.11) fails for every
      policy of every non-trivial task network.
  \end{itemize}
  We formalise the intended semantics, which is also what the proofs of Props.~2.12, 2.13 and
  2.15 use. A \emph{configuration} is a pair $(T',s)$ of executed tasks and current state
  (\lean{FMConfig}). The policy gives a \emph{valid instruction} $t$ at $(T',s)$
  (\lean{FMValid}) iff $\pi(T',s)=t$, $t\in T\setminus T'$, all $\prec$-predecessors of $t$ are
  in $T'$, and $\alpha(t)$ is applicable in $s$. A step (\lean{FMStep}) executes a valid
  instruction and moves to $(T'\cup\{t\},s')$ for any $s'\in\gamma(t,s)$. Then $\pi$ is a
  \emph{strong} OD solution iff every configuration reachable from $(\emptyset,\sI)$ is complete
  ($T'=T$) or has a valid instruction (\lean{StrongOD}), i.e.\ no run gets stuck; runs always
  terminate since each step executes a new task. It is a \emph{weak} OD solution iff some
  reachable configuration is complete (\lean{WeakOD}).
\end{deviation}

\begin{deviation}[Side condition of Def.~2.8]
  [CB21, Def.~2.8] restricts the policy itself: $\langle(T',s),t\rangle\in\pi$ only if
  $t\notin T'$ and no unexecuted task precedes $t$. We allow arbitrary partial functions
  \lean{FMPolicy := Finset I → Set F → Option I} and check the condition during execution
  (\lean{FMValid}); an instruction violating it makes the run stuck. Every policy in the sense of
  the paper is one in our sense, and every policy in our sense can be restricted to its valid
  instructions without changing its runs, so solution existence is unaffected.
\end{deviation}

\begin{deviation}[Belief states pointwise]
  [DB25] defines solvability of a \emph{belief state}, a set of pairs $(\tn,S)$, as solvability
  of all its members. We define solvability of a single pair (\lean{BeliefSolvable}) as an
  inductive predicate; solvability of a set is then the conjunction over its members. The
  recursion of [DB25] is well founded since $\tn\setminus t$ has fewer tasks, so the inductive
  reading is the intended one.
\end{deviation}

\section{Method-Based Policies}\label{sec:mp}

\begin{deviation}[Policies up to isomorphism]\label{dev:iso}
  A policy of [CB22, Def.~5] is a partial function on pairs $(\tn,s)$ that is not defined on two
  distinct isomorphic task networks for the same state; we formalise it as a partial function on
  concrete task networks (\lean{MPPolicy}) with that condition (\lean{MPPolicy.WellDefined}).
  Algorithm~1 looks up the instruction $(t,m)$ at an isomorphic entry $(\tn',s)$, but $t$ is then
  a task of $\tn'$, not of $\tn$. We make the implicit transport explicit: the instruction at
  $(\tn,s)$ is $(\sigma(t'),m)$ for an entry $\pi(\tn',s)=(t',m)$ with $t'\in\tn'$ and an
  isomorphism $\sigma:\tn'\to\tn$ (\lean{MPInstr}). If $\tn$ has non-trivial automorphisms,
  different $\sigma$ give different but automorphic tasks, whose successors are isomorphic.
\end{deviation}

\begin{deviation}[The quotient execution structure]\label{dev:quot}
  [CB22, Def.~6] first builds the execution structure on concrete pairs and then takes the
  quotient by isomorphism. We keep concrete nodes (\lean{MPNode}) and edges (\lean{MPStep}) and
  express the properties of the quotient $[L]$ directly:
  \begin{itemize}
    \item $[L]$ is finite iff the reachable nodes fall into finitely many isomorphism classes
      (\lean{MPFinite});
    \item $[L]$ is acyclic iff no non-empty path leads from a reachable node to a node isomorphic
      to it (\lean{MPAcyclic});
    \item terminal and goal nodes are invariant under isomorphism and are used as is.
  \end{itemize}
  Decomposition has a successor for \emph{every} fresh copy of the method's network; all of them
  are isomorphic and collapse in the quotient.
\end{deviation}

\begin{deviation}[Primitive tasks must be unconstrained]\label{dev:unc}
  [CB22, Def.~6] creates an edge for an instruction $(t,\varepsilon)$ with $t$ primitive without
  requiring that $t$ has no predecessors in $\tn$ (nor does Def.~5); the ordering of primitive
  tasks would then be meaningless. The progression algorithm of [CB22, Alg.~2] only executes
  tasks without predecessors. We require the executed task to be unconstrained
  (\lean{MPStep.prim}). Compound tasks may be decomposed regardless of their position, exactly
  as in Defs.~5--6.
\end{deviation}

\begin{deviation}[Ill-formed instructions]
  [CB22, Def.~5] demands $m=\varepsilon$ for primitive and $m=(\alpha(t),\tn_m)\in M$ for compound
  tasks. We allow any instruction; one violating this, or executing an inapplicable task, simply
  has no successors, so the node is terminal (and fails unless it is a goal node). This agrees
  with the paper's treatment of inapplicable tasks, for which $\gamma$ is undefined.
\end{deviation}

\begin{deviation}[Removing a task]
  [CB22] writes $\tn\setminus t$ for the restriction of $\tn$ to $T\setminus\{t\}$. We use
  progression \lean{HTN.TaskNetwork.progress} of \lean{HtnPlanning}, which removes $t$ and all
  pairs $(t,t')$ from $\prec$ but may keep pairs $(t',t)$. As $t$ is unconstrained no such pair
  with $t'\in T$ exists in a well-formed network, and in any case pairs involving tasks outside
  the network are ignored by every notion considered here.
\end{deviation}

\begin{deviation}[Strong cyclic solutions]\label{dev:sc}
  We formalise strong cyclic solutions exactly as [CB22, Def.~7] states them: every terminal
  node of $[L]$ is a goal node (\lean{IsMPStrongCyclicSolution}). This admits executions that
  never terminate, e.g.\ a policy that forever decomposes a recursive task, and differs from the
  usual notion of \textcite{cimatti:et:al:aij-03} (from every reachable node some goal node is
  reachable). The latter is provided for comparison as \lean{IsMPStrongCyclicFairSolution}. The
  literal reading is the one assumed by the membership algorithm for totally ordered problems in
  our notes, which uses triples $(c,s,\emptyset)$ for tasks that never finish.
\end{deviation}

\section{Proved Results}\label{sec:results}

\paragraph{Strong OD solutions are solvable belief states.}
For every primitive task network $\tn$, there is a policy that strongly executes $\tn$ in the
sense of Deviation~\ref{dev:od} iff $(\tn,\sI)$ is a solvable belief state
(\lean{FOND.Problem.strongOD_iff_beliefSolvable}). Hence strong OD plan existence of [CB21]
coincides with plan existence of [DB25]
(\lean{FOND.Problem.strongODSolvable_iff_fmSolvableDB}). This confirms that
Deviation~\ref{dev:od} captures the intended notion. The policy constructed from a solvable
belief state chooses, at each configuration, a task witnessing solvability of the remaining
network.

\printbibliography

\end{document}
